(1) 直接接触无对外功、无对外热，能量守恒给 \(C_1(T_{\rm eq}-T_1)+C_2(T_{\rm eq}-T_2)=0\)，因此 \(T_{\rm eq}=(C_1T_1+C_2T_2)/(C_1+C_2)\)。它是严格位于初温之间的加权算术平均，不能使用无权平均。
(2) 因为热容恒定，熵变为 \(\Delta S=C_1\ln(T_{\rm eq}/T_1)+C_2\ln(T_{\rm eq}/T_2)\)。对凹函数 \(\ln x\) 用 Jensen 不等式：\(\ln[(C_1T_1+C_2T_2)/C_\Sigma]\ge[C_1\ln T_1+C_2\ln T_2]/C_\Sigma\)，故 \(\Delta S\ge0\)；仅 \(T_1=T_2\) 时等号成立。
(3) 可逆热机不向外排热且过程熵守恒，故 \(C_\Sigma\ln T_*=C_1\ln T_1+C_2\ln T_2\)，即 \(T_*=T_1^{C_1/C_\Sigma}T_2^{C_2/C_\Sigma}\)。因外界只有功流，最大取功 \(W_{\max}=C_1T_1+C_2T_2-C_\Sigma T_*\)，其中 \(C_\Sigma=C_1+C_2\)。
(4) 加权算术—几何平均不等式给 \(T_*\le T_{\rm eq}\)，所以直接混合使可用功损失 \(C_\Sigma(T_{\rm eq}-T_*)\ge0\)。相同初态仅凭能量守恒允许多种末态；熵产生约束决定能否达到更低的最终温度以及可提取的机械功上界。
